% !TeX root = ../papers/nc-p3.tex
\section{二交換を仮定しない $z=1$ の混合補完}
\label{sec:klein-mixed-exchange-resolution}
節\ref{sec:klein-elliptic-completion}の規格化で $z=1$、
$\lambda\ne0,\pm1$ とし、隣接二切断の色別階数 $(2,2,1,1)$ を課す。
$P=\left(\begin{smallmatrix}0&1\\1&0\end{smallmatrix}\right)$ とする。
$X=X_\lambda,Y=Y_\lambda$ はともに可逆で $Y=\lambda PX$、$XP=PX$ である。
従って最初の二角 $A=A_0,B=A_1$ は必ず可逆であり、$T=YA^{-1}X,T'=YB^{-1}X$ を使える。
\begin{theorem}\label{thm:klein-mixed-exchange-resolution}
この繊維で標準正則となり得る全窓は、以下の表示とその二位置移動・文字置換で覆われる。
\[
 A=\begin{pmatrix}a&b\\b&d\end{pmatrix},\quad
 B=\begin{pmatrix}b&f\\g&b\end{pmatrix},\quad
 x=ad-fg\ne0,\quad t=\frac{a+d-(\lambda+1/\lambda)b}{ad-b^2}
   =\frac{(\lambda+1/\lambda)b-f-g}{b^2-fg}.
\]
この式と $\det A\det B\ne0$ は表示された側の二切断条件に必要十分である。
この側は共通倍率を除き有理既約四次元で、楕円曲線をパラメータに持つ五次元族の閉包に属する。
ここでの幾何学的一般点は標準Koszul左右AS正則、反対階数16、最小周期4、
指定始点の射影安定群は四元群である。印を忘れた像も四次元である。
\end{theorem}
\begin{proof}[全補完の分岐]
第二切断の二つの厚い行列 $N_0,N_1^{\mathsf t}$ の交差ブロックは全て $J=\boldsymbol1_{2\times2}$。
上の角を $C,D$ と書く。$C$ は $A$ の対角と $B$ の非対角、$D$ は $B$ の対角と
$A$ の非対角を使う。下の角はそれぞれ $T,T'$ の対応する成分を同じ順序で使う。
一つの厚い行列を $\binom{L}{R}(K\ S)$ と分解すると $LS=RK=J$ である。
上の角 $LK$ が特異なら $L$ の列が定数、または $K$ の行が定数である。
従ってその二行または二列が一致する。二列だけが一致する場合 $L$ は可逆なので
$S$ の二列も、下の角の二列も一致する。行についても同じである。

$C,D$ がともに特異な場合、片方の二列と他方の二行が一致すれば
$B=A^{\mathsf t}P$ または $B=PA^{\mathsf t}$ である。
それ以外では両方で二列だけ、または両方で二行だけが一致する。
前者で $B=AP,T'=PT$、しかも二つの下の角の両方で二列が一致する。
従って $T_{00}=T_{11},T_{01}=T_{10}$、$T$ と $A$ は $P$ と可換となり、
$C,D$ はともに $J$ の倍数で矛盾する。後者は転置で同じである。

残った二つの特異角表示を経路長4で排除する。
$B=A^{\mathsf t}P$ なら $T'=T^{\mathsf t}P$ であり、
$N_0$ の二つの列対は各々同一、$N_1^{\mathsf t}$ の二つの行対も各々同一である。
従って始点3の厚い二色の関係の係数空間はいずれも $(1,1,0,0),(0,0,1,1)$ の張る平面となる。
薄い二色は定数行である。始点2の長さ4で両端の二次関係を先に割ると、各総文字での次元は
$26,26,24,24$。中央の関係の挿入の関数体上の階数は $13,13,12,12$ で、
残る全次元は50である。具体的には端の各色の商を、厚いブロック $M$ の上の角を $G$、
左下を $H$ として $(-HG^{-1}\ I_2)$ で与え、薄い色では $(I_3\ {-\boldsymbol1})$ を使う。
左右のこの二射影に中央の関係の各項を通した行列を有理係数多項式環上で掃き出す。
\path{computations/verify_exchange_mixed_resolution.py} が独立な五変数で検算する。
階数6を保つ全特殊化でも次元は50以上であり、標準値35に反する。
$B=PA^{\mathsf t}$ は二位置移動で $A,B$ を転置し $\lambda$ を逆数にすれば同じである。

$C$ が可逆なら $N_0$ の下の角は $JC^{-1}J=tJ$。
$T_{00}=T_{11}$ から $A$ の非対角が等しく、$T'_{01}=T'_{10}$ から $B$ の対角が等しい。
これは $\lambda-1/\lambda\ne0$ による。もし $D$ も可逆なら下の角も定数で、
$T,T'$ は $P$ と可換、しかも $T'=PT$。従って $A,B$ も可換で $B=AP$、
$C$ が特異となり矛盾する。$D$ は対角が等しく非対角も等しい特異行列だから、
先の行または列の一致より $D=bJ$ である。$C$ の行列式は $x$ であり、$T_{00}=T'_{01}$ が
記載の $t$ の等式となる。$t(ad-fg)=a+d-f-g$ はその二式の和から従うので $tJ=JC^{-1}J$。
他方の下の角を $Z=\left(\begin{smallmatrix}T'_{00}&T_{01}\\T_{10}&T'_{11}\end{smallmatrix}\right)$ とすると、
残る条件は $b\det Z=Z_{00}+Z_{11}-Z_{01}-Z_{10}$ である。実際に両辺の差は
$T_{00}-T'_{01}$ の倍数で、商の分母は $\lambda\det A\det B$ を割る。
$Z$ も可逆である。さもなくば高位と低位の交換で上の角が二つとも特異となる。
この交換は $A,B$ を $T,T'$ にし $\lambda$ を逆数にするので、先の特異角分岐から
$T'=T^{\mathsf t}P$ または $PT^{\mathsf t}$、従って $B=A^{\mathsf t}P$ または $PA^{\mathsf t}$
となって $C$ が特異になる。ここで$Z$ の可逆性の議論に階数2の下界が必要な場合、
階数1なら $b\ne0,Z=J/b$ であり、二つの下の角が定数なので先と同じ矛盾である。
最後に $D$ が可逆な側は、二位置移動、奇文字への1の加法、全位置の $2,3$ の交換で
$(A,B,\lambda)$ を $(B,A,1/\lambda)$ として得る。これは始点の変更を明示した被覆である。
\end{proof}

\begin{proof}[特殊化と正則点]
$(tf-1)g=f+tb^2-(\lambda+1/\lambda)b$ と書ける。$tf-1$ で割った係数環は
$t$ を逆にした整域であり、その上で右辺は恒等的に零でない。
従って $g$ の一次方程式は原始的かつ既約で、$tf=1$ 上に別成分はない。
共通倍率を除けば四自由度である。

一般点を滑らかにする一変数を $\epsilon=z-1$ とし、節\ref{sec:klein-elliptic-completion}の
$A$ の左下を $c=b+\epsilon k$ とおく。$x=ad-fg$ を固定し、同節の
$f=f_0+f_1x,g=g_0+g_1x$ による $fg+x-ad=0$ を $k$ の式と見る。
全係数は $\epsilon=0$ で正則であり、$f_0=g_0=(a+d)/2$、
$f_1+g_1=-t$、$\partial f_1/\partial k=(1/\lambda-\lambda)/(2\det A)$ がその剰余である。
従って与えた $f,g$ を実現する剰余 $k$ があり、$f\ne g$ の一般点で方程式の
$k$ 微分は $(g-f)x(1/\lambda-\lambda)/(2\det A)\ne0$。形式的陰関数定理が持ち上げを与える。
同節の二次式を $p x^2+m x+n$, $q y^2+j y+o$ と書けば、正しい枝は
$y=-(2px+m+j)/(2q)$。$p,m,j$ は正則で、$e_1,h_1$ の単純極の主要部は逆、
$q=e_1h_1$ の極の位数は2である。従って $y=O(\epsilon^2)$ で $e,h$ は $b,b$ に戻る。
一般の剰余で $o=(b-c)(z^2b-c)/(z+1)^2$ の主要係数は非零なので一般繊維では $y\ne0$。
以上は五次元族からの特殊化であり、既約性により表示全体が閉包に入る。

二つの有理点 $(a,b,d,f,\lambda,g)=(2,3,7,5,2,2/5),(3,5,7,4,3,87/17)$ では
反対切断が可逆で、全始点の次数5までの次元は標準値である。文字別の四次中心次元は
$(3,2,0,0)$。前掲の検算が有理数上で四式を選び、その商の素数1009での次元を全始点で
次数13まで調べると $(1+t+t^2+t^3)^4$ の係数である。
五次元族の一般点から全次数の標準Hilbert関数が下界であるため、
命題\ref{prop:klein-rational-full-opposite}の自由加群・Koszul双対の証明が適用される。
従って両点で全次数の正則性が従い、有限Frobenius条件と各内部次数のbar複体から
本表示の幾何学的一般点も正則である。

第一の有理点を素数1753で還元すると、各位置で被約二十点、三点直線なし、四点平面52枚で
五点平面はない。各点を通る四点平面の数を保存するように枠の像を全列挙した
射影写像数の行は $(4,0,0,4),(0,4,4,0),(0,4,4,0),(4,0,0,4)$ である。
\path{computations/verify_exchange_mixed_configurations.py} は各配置を独立にも消去して照合する。
有限エタールな配置の固有な枠比較よりこの計算は上界を与え、位置0と2は標数零および
一般点で同値にならない。よって最小周期は4である。各配置の射影群は指定四元群を含むため
位数4である。窓を保つ四位置の写像は位置0の写像で決まる。実際二十点間の一意な点対応が
残りの配置の各点の像を決め、各配置は射影枠を含む。従って安定群は四元群であり、印を忘れる
ファイバーも有限である。厚い角の共通倍率以外の連続な対角残余がないことは、薄い
48語の文字指数行列の階数12と、上の角の文字指数差がその整数行格子に入ることから従う。
\end{proof}
